\section{Introdução}

Este trabalho analisa um sistema linear e invariante no tempo (LIT) discreto,
excitado por processos estocásticos. Os dados foram gerados a partir desse sistema
e organizados em quatro conjuntos. O objetivo é caracterizar os sinais de entrada e
de saída nos domínios do tempo e da frequência, estimar a relação entre entrada e
saída, identificar um modelo para o sistema e verificar até que ponto esse modelo
descreve os dados.

A análise é dividida em três partes, que acompanham a organização proposta no
enunciado. A primeira parte trata do conjunto com entrada aproximadamente branca e
inclui a identificação e a validação do modelo. A segunda parte repete a análise
temporal e espectral para um conjunto com entrada colorida e compara os dois casos.
A terceira parte aplica o modelo identificado a um conjunto de teste para verificar
a presença de anomalias por meio dos resíduos.

Ao longo do texto, procura-se não apenas apresentar os gráficos e as estimativas,
mas também comentar como a autocorrelação, a densidade espectral de potência, o
espectro cruzado, a coerência e os resíduos ajudam a avaliar a qualidade e a
confiabilidade do modelo. Os resultados numéricos e as figuras foram produzidos por
um conjunto de rotinas em Python, descritas de forma resumida na próxima seção, de
modo que podem ser reproduzidos.
